\item Una molécula triatómica puede tener una configuración lineal, como el $\text{CO}_2$ (figura P8.40a), o puede ser no lineal, como el $\text{H}_2\text{O}$ (figura P8.40b). Suponga que la temperatura de un gas de moléculas triatómicas es suficientemente baja como para que el movimiento vibratorio sea despreciable. ¿Cuál es el calor específico molar a volumen constante, expresado como múltiplo de la constante universal de los gases, (a) si las moléculas son lineales y (b) si las moléculas son no lineales? En altas temperaturas, una molécula triatómica tiene dos modos de vibración, y cada uno aporta $\frac{1}{2}R$ al calor específico molar para su energía cinética y otro $\frac{1}{2}R$ para su energía potencial. Identifique el calor específico molar de alta temperatura a volumen constante para un gas ideal triatómico de (c) moléculas lineales y (d) moléculas no lineales. (e) Explique cómo se pueden emplear los datos de calor específico para determinar si una molécula triatómica es lineal o no lineal. ¿Los datos en la tabla 8.2 son suficientes para hacer esta determinación?
